\documentclass[portrait,a0,final]{a0poster}

%nieco od Jakuba:
  \pdfpagewidth=\paperwidth                           % zabezpecuje, aby to v~pdf malo rovnaky rozmer ako v~dvi,
  \pdfpageheight=\paperheight                         % tie za beznych okolnosti nemusia moc dobre korelovat

%Slovencina
%\usepackage{slovak}

%\usepackage{slovak}
%\usepackage[utf8]{inputenc}

%\usepackage[slovak]{babel}
%\usepackage[cp1250]{inputenc}
%\usepackage[T1]{fontenc} %pekne makcene

%male popisy obrazkov a~grafika
\usepackage[font=small,margin=0.5cm]{caption} % margin reguluje okraje popisu obrazka (v pripade, ze je na sirku strany a~ma viac ako 1 riadok)
\usepackage[dvips]{graphicx}
\usepackage{wrapfig}
\usepackage[usenames,dvipsnames]{color}
\usepackage{epstopdf}   %bez tohto pdflatex nezoberie eps obazky
\usepackage{listings}
\lstset{ 
	basicstyle=\footnotesize,
    tabsize=3,
} 

%grid obrazkov
\usepackage{subfig}

\newcounter{eqn}
\renewcommand*{\theeqn}{\alph{eqn})}
\newcommand{\num}{\refstepcounter{eqn}\text{\theeqn}\;}

\makeatletter
\newcommand{\putindeepbox}[2][0.7\baselineskip]{{%
    \setbox0=\hbox{#2}%
    \setbox0=\vbox{\noindent\hsize=\wd0\unhbox0}
    \@tempdima=\dp0
    \advance\@tempdima by \ht0
    \advance\@tempdima by -#1\relax
    \dp0=\@tempdima
    \ht0=#1\relax
    \box0
}}
\makeatother

%farebne tabulky
\usepackage{colortbl}
\usepackage[table]{xcolor}
%na otacanie tabuliek
\usepackage{rotating}

%odsadenie prveho odstavca
\usepackage{indentfirst}

%Matematicke vyrazy
\usepackage{amsfonts}
\usepackage{amsmath}
\usepackage{amssymb}

\usepackage{verbatim}
\usepackage[official]{eurosym}
\usepackage{url}

%algoritmy
\usepackage[lined,boxed]{algorithm2e}

%definicie
\newcommand{\p}{\partial}
 \def\epsilon{\varepsilon}
 \def\Bf#1{\mathbf{#1}}

%TRUBAC
\usepackage[left=2cm,right=2cm,top=0cm,bottom=0cm, paperwidth=841mm, paperheight=1201.42mm]{geometry}
\usepackage{a0poster}
\usepackage{multicol}
\usepackage[]{indentfirst,latexsym}

\usepackage{pstricks,pst-grad}

\renewcommand{\rmdefault}{phv} % Arial
\renewcommand{\sfdefault}{phv} % Arial

% Hrubka oddelovaca medzi stlpcami
\columnseprule=2pt
% Medzera medzi oddelovacmi stlpcov a textom
\columnsep=100pt

%====================================================================================================================================================
%====================================================================================================================================================

\begin{document}

% Font + riadkovanie
\renewcommand{\baselinestretch}{1.5}
\fontsize{25pt}{20pt}
\selectfont

\begin{center}
\includegraphics[width=\textwidth]{header.eps}
\end{center}

\centerline{\huge \bf{Cayley graphs of given diameter or girth on linear groups}}
\vskip 1cm 
\centerline{\Large Mat\'u\v{s} Behun}
\vskip 0.5cm
\leaders\vrule width \textwidth\vskip 3pt % or other desired thickness
%\noindent\hrulefill

\begin{multicols}{3}
\paragraph{Motivation}
\begin{itemize}
	\item In its simplest form, networks can be modeled by graphs with nodes as vertices and links between them as edges.
	\item In modeling interconnection networks by graphs one often considers these restrictions:
	\begin{itemize}
		\item limitation on the number of physical links leaving a node
		\item limitation on accessibility, that is, any two nodes should be accessible using at most a certain number of physical links.
		\item limitation on length of smallest cycle in graph.
	\end{itemize}
\end{itemize}

In graph theory this leads to two important problems: the {\em degree/diameter problem} to construct the largest possible graphs of a given maximum degree and a given diameter, and the {\em degree/girth problem} to construct the smallest possible regular graphs of a given degree and a given girth; in both cases the adjectives `large' and `small' refer to the order (i.e., the number of vertices) of a graph.
\paragraph{Degree/diameter problem}
~\\
There is a theoretical upper bound on the largest order of a graph of maximum degree $d\ge 2$ and diameter $k\ge 1$.
\begin{equation}\label{eq:Moore}
	\begin{split}
		n_{d,k} \leq M_{d,k}    & = 1 + d + d(d - 1) + \dots + d(d - 1)^{k-1}  \\
								& = 1 + d(1 + (d - 1) + \dots + (d - 1)^{k-1}) \\
								& = \begin{cases}
								1+d\frac{(d-1)^{k}-1}{d-2}, & \text{if}\ d > 2 \\
								2k+1, & \text{if}\ d=2
									\end{cases}
	\end{split}
\end{equation} ~\\~\\
Graphs of maximum degree $d$ and diameter $k$ with $M_{d,k}$ vertices are called {\em Moore graphs}. Moore graphs are rare; they exist only for the following degrees and diameters: 
~\\
\begin{itemize}
	\item If $d = 2$ for any $k \geq 1$
	\item If $k = 1$ for any $d \geq 2$
	\item If $k = 2$ and $d \in \{3, 7 \}$, and possibly $57$
\end{itemize} ~\\
In the remaining cases one tries to construct graphs of maximum degree $d$ and diameter $k$ of order as close to the Moore bound $M_{d,k}$ as possible.

\begin{figure}
	\centering
	\includegraphics[scale=0.7]{Hoffman-Singleton_graph.png}

	\textbf{Hoffman-Singleton graph with $k=2$ and $d=7$}
\end{figure}

~\\
\paragraph{Construction of large graphs}
~\\
Finding $(d,k)$-graph of large order is approached by many techniques, mostly using combinatorics on words or various algebraic structures. \\~\\
We present two examples of bounds arising from such constructions: ~\\
Using combinatorics on words, Baskoro and Miller(1993) \cite{Bas-Mil} proved that 
\begin{equation*}
	n_{d,k} \geq \left( \frac{d}{2} \right)^{k} +  \left( \frac{d}{2} \right)^{k-1}     
\end{equation*}
With the help of finite fields Bevan, Erskine and Lewis(2017) \cite{Bev-Ers} observed that modified Brown graphs give the bound
\begin{equation*}
	n_{d,2} \geq d^{2} - 2d^{1+\varepsilon}     
\end{equation*}
where $\varepsilon$ depends on results about graphs between consecutive primes. 

\paragraph{Cayley graphs}  
~\\
Let $\Gamma$ be a group and let $S\subset \Gamma$ be a symmetric unit-free generating set for $\Gamma$; that is, we require that $S=S^{-1}$ and $1\notin S$. 
The $\textit{Cayley graph}$ $C(\Gamma,S)$ is the graph with vertex set $\Gamma$ in which vertices $a,b$ are adjacent if $a^{-1}b\in S$. ~\\

By $C_{d,k}$ we denote the largest order of Cayley graph of degree $d$ and diameter $k$. ~\\
The best currently known lower bound on diameter $2$ for an infinite family of degrees is due to \v{S}iagiov\'a and \v{S}ir\'a\v{n}(2012): ~\\
Let $D = \{ 2^{2m+\mu}+    (2+\delta)2^{m+1}-6,m \geq 1, \mu \in \{0,1\} \}$. Then, for every $d\in D$ one has $C_{d,2} > d^{2} - 6\sqrt{2}d^{3/2}$.

\paragraph{Special linear group}  
~\\
%\begin{definition}[General and special linear group] 
Let $q$ be a power of a prime and let $GF(q)$ be the Galois field of order $q$. The {\em general linear group} $GL(m,q)$ consists of all non-singular $m\times m$ matrices over $GF(q)$ under multiplication. \\~\\
The special linear group is the subgroup of $GL(m,q)$ consiting of matrices with determinant $1$.
~\\
$|GL(m,q)| = (q^m - 1)(q^m - q) \cdots (q^m - q^{m-1})$
~\\
$|SL(m,q)| = |GL(m,q)|/(q-1)$
~\\
%\end{theorem}
\paragraph{Computer search}
~\\
We created a program generating Cayley graphs of given degree and diameter based on $SL(2,q)$. \\~\\
As an example, consider the problem of generating a Cayley graph for the group $SL(2,5)$ of diameter $2$ and minimum possible degree.
~\\
In order to find graphs with given $d$ and $k$ we generted Cayley graphs from random symmetric unit-free sets $S$ of $SL(2,5)$.
The order of $SL(2,5)$ is $120$, so that by the Moore bound the first feasible degree (equal to sthe size of a generating set) would be $12$.

The group $SL(2,5)$ has only one involution; all the remaining $118$ elements form $59$ pairs of the form $\{x,x^{-1}\}$. \\~\\
To check all Cayley graphs $C(G,S)$ for $G=SL(2,5)$ and $|S|=12$ we would have to generate all the ${59 \choose 6} = 45057474$ possibilities for $S$ and then check for the diameter of the resulting Cayley graphs. \\~\\
For the group $G=SL(2,5)$ and generating sets $S$ such that $|S|=d$ by our randomized algorithm we found the following number of generating sets giving Cayley graphs $C(G,S)$ of diameter $2$: \\~\\

\begin{tabular}[htbp]{l*{10}{c}r}
		$d$ & $13$ & $14$ & $15$ & $16$ & $17$ & $18$ & $19$ & $20$ & $21$ \\ 
		\hline
		Found graphs & $0$ & $0$ & $1$ & $107$ & $345$ & $2451$ & $4120$ & $11669$ & $14926$ \\
\end{tabular}

\begin{figure}
	\centering
	\includegraphics[scale=0.40]{example.png} 

	\textbf{Example of graph generated on $SL(2,3)$}
\end{figure}
~\\~\\~\\~\\
\paragraph{Genearation of Cayley graph}
\lstinputlisting[language=Perl]{cayley_graph_generation.pl}

\paragraph{Check of diameter of Cayley graph}
\lstinputlisting[language=Perl]{check_diameter.pl}


\paragraph{Degree/girth problem}
~\\
Finding regular graphs with smallest possible order denoted $n_{d,g}$ with given degree $d$ and girth $g\geq3$ is known as {\em degree/girth problem}. Motivation for finding such graphs could arise from constructing graphs with no cycles of length less than $g$ and its similarity to {\em degree/diameter problem}.
\begin{equation*}
	n_{d,g} \geq M(d,g) = 1 + d + d(d - 1) + d(d - 1)^{2} + \dots + d(d - 1)^{\frac{g-3}{2}}
\end{equation*}
and for even $g$
\begin{equation*}
	n_{d,g} \geq M(d,g) = 2(1 + (d - 1) + (d - 1)^{2} + \dots + (d - 1)^{\frac{g-2}{2}})
\end{equation*}

For odd $g$ the Moore bound is obtained in the same way as the Moore bound for the degree/diameter problem for the diameter $k=(g-1)/2$. For even $g$ we start with an edge and from its incident vertices we continue growing two spanning trees until we reach depth $(g-2)/2$.

Graphs of degree $d$ and girth $g$ with order equal to $M(d,g)$ are called Moore graphs or cages. We list all cases with $n_{d,g} = M_{d,g}$ below:

\begin{itemize}
	\item{\makebox[5cm]{For $d=2$,\hfill}$g \geq 2$ - circles}
    \item{\makebox[5cm]{For $g=3$,\hfill}$d \geq 2$ - complete graphs}
    \item{\makebox[5cm]{For $g=4$,\hfill}$d \geq 2$ - complete bipartite graphs}
    \item{\makebox[5cm]{For $g=5$,\hfill}$d=2$ - circle of length $5$}
    \item[]{\makebox[5cm]{\hfill}$d=3$ - Peterssen graph}
    \item[]{\makebox[5cm]{\hfill}$d=7$ - Hoffman-Singleton graph}
    \item[]{\makebox[5cm]{\hfill}$d=57$ - this value has not been excluded but no such graph has been found yet }
    \item{\makebox[7cm]{For $g \in {6,8,12}$,\hfill}if $d-1$ is a prime power}
\end{itemize}
\begin{figure}
	\centering
	\includegraphics[scale=0.4]{Heawood_graph.png} 

	\textbf{Heawood graph with $d=3$ and $g=6$}
\end{figure}
\end{multicols}

\vfill
\begin{center}
\makebox[\textwidth]{\includegraphics[width=\paperwidth]{footer/footer.pdf}}
\end{center}
\end{document}
